\chapter{RE-007 --- Migrate Framed Shell Onto Runtime Engine Seams}\label{record-155}

\noindent{\footnotesize\color{BookMuted}\path{docs/design/RE-007-migrate-framed-shell-onto-runtime-engine-seams.md}\par}\smallskip

\emph{Cycle for proving the runtime engine through the framed shell instead of leaving the most visible shell state on ad hoc branch structure}

Legend:

\begin{itemize}
\tightlist
\item
  \href{https://github.com/flyingrobots/bijou/blob/1853ad902f6bf219326409e7f743a5a1d6517dbc/docs/legends/RE-runtime-engine.md}{RE --- Runtime Engine}
\end{itemize}

Depends on:

\begin{itemize}
\tightlist
\item
  \href{https://github.com/flyingrobots/bijou/blob/cbd70612140e7a450d764b3c7be1c6af878b99bc/docs/design/RE-001-define-runtime-engine-architecture.md}{RE-001 --- Define Runtime Engine Architecture}
\item
  \href{https://github.com/flyingrobots/bijou/blob/cbd70612140e7a450d764b3c7be1c6af878b99bc/docs/design/RE-002-promote-first-class-state-machine-and-view-stack.md}{RE-002 --- Promote First-Class State Machine and View Stack}
\item
  \href{https://github.com/flyingrobots/bijou/blob/cbd70612140e7a450d764b3c7be1c6af878b99bc/docs/design/RE-003-retain-layout-trees-and-layout-invalidation.md}{RE-003 --- Retain Layout Trees and Layout Invalidation}
\item
  \href{https://github.com/flyingrobots/bijou/blob/cbd70612140e7a450d764b3c7be1c6af878b99bc/docs/design/RE-004-route-input-through-layouts-and-layer-bubbling.md}{RE-004 --- Route Input Through Layouts and Layer Bubbling}
\item
  \href{https://github.com/flyingrobots/bijou/blob/cbd70612140e7a450d764b3c7be1c6af878b99bc/docs/design/RE-005-buffer-commands-and-effects-separately.md}{RE-005 --- Buffer Commands and Effects Separately}
\item
  \href{https://github.com/flyingrobots/bijou/blob/cbd70612140e7a450d764b3c7be1c6af878b99bc/docs/design/RE-006-formalize-component-layout-and-interaction-contracts.md}{RE-006 --- Formalize Component Layout and Interaction Contracts}
\end{itemize}

\section{Sponsor human}\label{record-155-sponsor-human}

A maintainer evolving the framed shell who needs to explain layer ownership, routing, and command dispatch in runtime-engine terms instead of tracing ad hoc branch order.

\section{Sponsor agent}\label{record-155-sponsor-agent}

An agent inspecting or testing the framed shell that needs to recover runtime-backed view stacks, retained layout geometry, and buffered command/effect truth through the same vocabulary the runtime engine exposes everywhere else.

\section{Non-goals}\label{record-155-non-goals}

\begin{itemize}
\tightlist
\item
  Rewriting shell rendering or visual output. This cycle migrates ownership and routing, not painting.
\item
  Changing outward shell behavior. The shell should act the same after each slice; only the internal truth source changes.
\item
  Migrating non-shell runtime consumers. Other runtime adopters are separate cycles.
\end{itemize}

\section{Accessibility / assistive reading posture}\label{record-155-accessibility-assistive-reading-posture}

Not directly affected. This cycle changes internal ownership, not user-facing output. Assistive reading contracts remain wherever the shell already honors them.

\section{Localization / directionality posture}\label{record-155-localization-directionality-posture}

Not directly affected. Shell localization surfaces are unchanged by this migration. Directionality remains the responsibility of the rendering layer, which this cycle does not touch.

\section{Agent inspectability / explainability posture}\label{record-155-agent-inspectability-explainability-posture}

Central to this cycle. Every slice must leave the shell's runtime state inspectable through the same \texttt{R\allowbreak{}u\allowbreak{}n\allowbreak{}t\allowbreak{}i\allowbreak{}m\allowbreak{}e\allowbreak{}V\allowbreak{}i\allowbreak{}e\allowbreak{}w\allowbreak{}S\allowbreak{}t\allowbreak{}a\allowbreak{}c\allowbreak{}k}, \texttt{R\allowbreak{}u\allowbreak{}n\allowbreak{}t\allowbreak{}i\allowbreak{}m\allowbreak{}e\allowbreak{}R\allowbreak{}e\allowbreak{}t\allowbreak{}a\allowbreak{}i\allowbreak{}n\allowbreak{}e\allowbreak{}d\allowbreak{}L\allowbreak{}a\allowbreak{}y\allowbreak{}o\allowbreak{}u\allowbreak{}t\allowbreak{}s}, and runtime buffer vocabulary that agents already use for non-shell runtime work. The agent hill is the primary proof surface.

\section{Why this cycle exists}\label{record-155-why-this-cycle-exists}

The runtime engine now has explicit seams for:

\begin{itemize}
\tightlist
\item
  state machines
\item
  view stacks
\item
  retained layouts
\item
  layout-driven input routing
\item
  buffered commands and effects
\end{itemize}

But the framed shell still carries its own most visible runtime truths in older shell-local structures. The shell can already describe a layer stack, but that stack is still reconstructed as a plain descriptor array rather than living as a first-class runtime object.

This cycle exists to stop leaving the highest-visibility shell behavior outside the runtime engine. The shell should become the proving surface for the engine seams rather than the place where those seams get bypassed.

\section{Human users / jobs / hills}\label{record-155-human-users-jobs-hills}

\subsection{Primary human users}\label{record-155-primary-human-users}

\begin{itemize}
\tightlist
\item
  maintainers evolving the framed shell
\item
  builders inspecting layered shell behavior
\item
  contributors debugging shell routing, footer truth, and overlay ownership
\end{itemize}

\subsection{Human jobs}\label{record-155-human-jobs}

\begin{enumerate}
\def\labelenumi{\arabic{enumi}.}
\tightlist
\item
  Inspect the shell's visible layers as a real runtime-backed stack.
\item
  Explain which shell layer is active and which layer sits underneath it without reverse-engineering branch order.
\item
  Move shell routing and command/effect dispatch onto the runtime seams in later slices without re-inventing shell-local state models.
\end{enumerate}

\subsection{Human hill}\label{record-155-human-hill}

A maintainer can point at the framed shell and explain its active layer stack in runtime-engine terms instead of saying ``the shell works because these branches happen in this order.''

\section{Agent users / jobs / hills}\label{record-155-agent-users-jobs-hills}

\subsection{Primary agent users}\label{record-155-primary-agent-users}

\begin{itemize}
\tightlist
\item
  agents reading shell state to explain focus or overlay ownership
\item
  agents writing tests around shell layering
\item
  agents preparing later shell migration slices for routing and buffered shell work
\end{itemize}

\subsection{Agent jobs}\label{record-155-agent-jobs}

\begin{enumerate}
\def\labelenumi{\arabic{enumi}.}
\tightlist
\item
  Read the framed shell's view stack through the same \texttt{R\allowbreak{}u\allowbreak{}n\allowbreak{}t\allowbreak{}i\allowbreak{}m\allowbreak{}e\allowbreak{}V\allowbreak{}i\allowbreak{}e\allowbreak{}w\allowbreak{}S\allowbreak{}t\allowbreak{}a\allowbreak{}c\allowbreak{}k} vocabulary used elsewhere in the runtime engine.
\item
  Distinguish the workspace root from dismissible overlays without re-deriving shell semantics.
\item
  Prepare later migration slices for retained-layout-driven shell routing and runtime-buffer-backed shell command/effect dispatch.
\end{enumerate}

\subsection{Agent hill}\label{record-155-agent-hill}

An agent can inspect the framed shell and recover a root-based runtime view stack whose active layer, underlying layer, and blocking semantics match the shell's real behavior.

\section{Human playback}\label{record-155-human-playback}

\begin{enumerate}
\def\labelenumi{\arabic{enumi}.}
\tightlist
\item
  A workspace is open in a framed app.
\item
  The workspace is the root runtime view.
\item
  The user opens settings.
\item
  The shell pushes a dismissible settings overlay onto the runtime view stack.
\item
  The user opens search on top of settings.
\item
  The shell pushes a dismissible search overlay onto the same runtime view stack.
\item
  The active runtime view now truthfully describes the search layer, while the underlying view is still settings.
\item
  Footer hints and help surfaces can inspect that same stack instead of rediscovering shell state from separate branches.
\end{enumerate}

\section{Agent playback}\label{record-155-agent-playback}

\begin{enumerate}
\def\labelenumi{\arabic{enumi}.}
\tightlist
\item
  An agent inspects a framed shell model with settings and search open.
\item
  The agent asks for the shell's runtime-backed view stack.
\item
  The stack returns:

  \begin{itemize}
  \tightlist
  \item
    a non-dismissible workspace root
  \item
    a dismissible settings overlay
  \item
    a dismissible search overlay
  \end{itemize}
\item
  The active runtime view reports the same metadata that shell help/footer introspection already uses.
\item
  The agent can now reason about shell overlay ownership with the same runtime vocabulary used by retained layouts and buffered runtime work.
\end{enumerate}

\section{Linked invariants}\label{record-155-linked-invariants}

\begin{itemize}
\tightlist
\item
  \href{https://github.com/flyingrobots/bijou/blob/cbd70612140e7a450d764b3c7be1c6af878b99bc/docs/invariants/runtime-truth-wins.md}{Runtime Truth Wins}
\item
  \href{https://github.com/flyingrobots/bijou/blob/cbd70612140e7a450d764b3c7be1c6af878b99bc/docs/invariants/state-machine-and-view-stack-are-distinct.md}{State Machine and View Stack Are Distinct}
\item
  \href{https://github.com/flyingrobots/bijou/blob/cbd70612140e7a450d764b3c7be1c6af878b99bc/docs/invariants/topmost-layer-dismisses-first.md}{Topmost Layer Dismisses First}
\item
  \href{https://github.com/flyingrobots/bijou/blob/cbd70612140e7a450d764b3c7be1c6af878b99bc/docs/invariants/layout-owns-interaction-geometry.md}{Layout Owns Interaction Geometry}
\item
  \href{https://github.com/flyingrobots/bijou/blob/cbd70612140e7a450d764b3c7be1c6af878b99bc/docs/invariants/commands-change-state-effects-do-not.md}{Commands Change State, Effects Do Not}
\item
  \href{https://github.com/flyingrobots/bijou/blob/cbd70612140e7a450d764b3c7be1c6af878b99bc/docs/invariants/tests-are-the-spec.md}{Tests Are the Spec}
\end{itemize}

\section{First slice in this cycle}\label{record-155-first-slice-in-this-cycle}

The first honest slice is not ``rewrite the whole shell.'' It is:

\begin{itemize}
\tightlist
\item
  make the framed shell layer stack a real \texttt{R\allowbreak{}u\allowbreak{}n\allowbreak{}t\allowbreak{}i\allowbreak{}m\allowbreak{}e\allowbreak{}V\allowbreak{}i\allowbreak{}e\allowbreak{}w\allowbreak{}S\allowbreak{}t\allowbreak{}a\allowbreak{}c\allowbreak{}k}
\item
  keep the existing \texttt{d\allowbreak{}e\allowbreak{}s\allowbreak{}c\allowbreak{}r\allowbreak{}i\allowbreak{}b\allowbreak{}e\allowbreak{}F\allowbreak{}r\allowbreak{}a\allowbreak{}m\allowbreak{}e\allowbreak{}L\allowbreak{}a\allowbreak{}y\allowbreak{}e\allowbreak{}r\allowbreak{}S\allowbreak{}t\allowbreak{}a\allowbreak{}c\allowbreak{}k\allowbreak{}(\allowbreak{}.\allowbreak{}.\allowbreak{}.\allowbreak{})}, \texttt{a\allowbreak{}c\allowbreak{}t\allowbreak{}i\allowbreak{}v\allowbreak{}e\allowbreak{}F\allowbreak{}r\allowbreak{}a\allowbreak{}m\allowbreak{}e\allowbreak{}L\allowbreak{}a\allowbreak{}y\allowbreak{}e\allowbreak{}r\allowbreak{}(\allowbreak{}.\allowbreak{}.\allowbreak{}.\allowbreak{})}, and \texttt{u\allowbreak{}n\allowbreak{}d\allowbreak{}e\allowbreak{}r\allowbreak{}l\allowbreak{}y\allowbreak{}i\allowbreak{}n\allowbreak{}g\allowbreak{}F\allowbreak{}r\allowbreak{}a\allowbreak{}m\allowbreak{}e\allowbreak{}L\allowbreak{}a\allowbreak{}y\allowbreak{}e\allowbreak{}r\allowbreak{}(\allowbreak{}.\allowbreak{}.\allowbreak{}.\allowbreak{})} helpers as compatibility surfaces
\item
  expose a runtime-backed \texttt{d\allowbreak{}e\allowbreak{}s\allowbreak{}c\allowbreak{}r\allowbreak{}i\allowbreak{}b\allowbreak{}e\allowbreak{}F\allowbreak{}r\allowbreak{}a\allowbreak{}m\allowbreak{}e\allowbreak{}R\allowbreak{}u\allowbreak{}n\allowbreak{}t\allowbreak{}i\allowbreak{}m\allowbreak{}e\allowbreak{}V\allowbreak{}i\allowbreak{}e\allowbreak{}w\allowbreak{}S\allowbreak{}t\allowbreak{}a\allowbreak{}c\allowbreak{}k\allowbreak{}(\allowbreak{}.\allowbreak{}.\allowbreak{}.\allowbreak{})} helper so tests, humans, and agents can inspect the shell directly in runtime terms
\end{itemize}

This slice proves that the shell can adopt runtime-backed view ownership without changing outward shell behavior first.

\section{Implementation outline}\label{record-155-implementation-outline}

\begin{enumerate}
\def\labelenumi{\arabic{enumi}.}
\tightlist
\item
  Promote shell layer descriptors into runtime-backed view layers with stable ids, blocking semantics, dismissibility, and attached descriptor models.
\item
  Make the workspace the root runtime view.
\item
  Push page-modal, settings, help, notification-center, search/command-palette, and quit-confirm overlays onto the runtime stack in the shell's existing order.
\item
  Derive the legacy descriptor helpers from that runtime view stack instead of rebuilding plain arrays separately.
\item
  Export the runtime-backed helper through the public framed-app API and the package index.
\end{enumerate}

\section{Later slices still in this cycle}\label{record-155-later-slices-still-in-this-cycle}

\begin{itemize}
\tightlist
\item
  retained-layout-driven shell routing
\item
  runtime-buffer-backed shell command/effect dispatch
\item
  removing remaining shell-local branch structures that duplicate runtime ownership
\end{itemize}

\section{Tests to write first}\label{record-155-tests-to-write-first}

\begin{itemize}
\tightlist
\item
  cycle test proving \texttt{R\allowbreak{}E\allowbreak{}-\allowbreak{}0\allowbreak{}0\allowbreak{}7} is now the active runtime-engine cycle
\item
  app-frame regression proving the framed shell layer stack is now backed by \texttt{R\allowbreak{}u\allowbreak{}n\allowbreak{}t\allowbreak{}i\allowbreak{}m\allowbreak{}e\allowbreak{}V\allowbreak{}i\allowbreak{}e\allowbreak{}w\allowbreak{}S\allowbreak{}t\allowbreak{}a\allowbreak{}c\allowbreak{}k}
\item
  export regression proving the new helper is public through \texttt{@\allowbreak{}f\allowbreak{}l\allowbreak{}y\allowbreak{}i\allowbreak{}n\allowbreak{}g\allowbreak{}r\allowbreak{}o\allowbreak{}b\allowbreak{}o\allowbreak{}t\allowbreak{}s\allowbreak{}/\allowbreak{}b\allowbreak{}i\allowbreak{}j\allowbreak{}o\allowbreak{}u\allowbreak{}-\allowbreak{}t\allowbreak{}u\allowbreak{}i}
\item
  signpost regression proving \texttt{P\allowbreak{}L\allowbreak{}A\allowbreak{}N\allowbreak{}.\allowbreak{}m\allowbreak{}d}, \texttt{B\allowbreak{}E\allowbreak{}A\allowbreak{}R\allowbreak{}I\allowbreak{}N\allowbreak{}G\allowbreak{}.\allowbreak{}m\allowbreak{}d}, and the RE legend stop treating \texttt{R\allowbreak{}E\allowbreak{}-\allowbreak{}0\allowbreak{}0\allowbreak{}7} like a backlog stub
\end{itemize}

\section{Risks / unknowns}\label{record-155-risks-unknowns}

\begin{itemize}
\tightlist
\item
  the shell still mixes routing, layout, and command/effect concerns, so this first slice must stay narrow instead of pretending the whole migration is done
\item
  keeping the compatibility descriptor helpers is necessary now, but later slices should avoid duplicating runtime truth unnecessarily
\item
  retained-layout-driven shell routing may want richer view metadata than the first runtime-backed stack carries here
\end{itemize}

\section{Retrospective}\label{record-155-retrospective}

What this first slice lands:

\begin{itemize}
\tightlist
\item
  a runtime-backed shell layer stack
\item
  a public \texttt{d\allowbreak{}e\allowbreak{}s\allowbreak{}c\allowbreak{}r\allowbreak{}i\allowbreak{}b\allowbreak{}e\allowbreak{}F\allowbreak{}r\allowbreak{}a\allowbreak{}m\allowbreak{}e\allowbreak{}R\allowbreak{}u\allowbreak{}n\allowbreak{}t\allowbreak{}i\allowbreak{}m\allowbreak{}e\allowbreak{}V\allowbreak{}i\allowbreak{}e\allowbreak{}w\allowbreak{}S\allowbreak{}t\allowbreak{}a\allowbreak{}c\allowbreak{}k\allowbreak{}(\allowbreak{}.\allowbreak{}.\allowbreak{}.\allowbreak{})} helper
\item
  compatibility descriptor helpers derived from runtime truth
\item
  signposts that say \texttt{R\allowbreak{}E\allowbreak{}-\allowbreak{}0\allowbreak{}0\allowbreak{}7} is active instead of merely queued
\end{itemize}

What this first slice does not land:

\begin{itemize}
\tightlist
\item
  no retained-layout-driven shell routing yet
\item
  no runtime-buffer-backed shell command/effect dispatch yet
\item
  no removal of every shell-local branch
\end{itemize}

What the second slice lands:

\begin{itemize}
\tightlist
\item
  route key ownership through the runtime view stack instead of choosing shell owners from ad hoc top-layer branches
\item
  use retained shell drawer layouts so pointer ownership for settings and notification center comes from runtime hit-testing instead of drawer-specific inside/outside helpers alone
\end{itemize}

What the third slice lands:

\begin{itemize}
\tightlist
\item
  a workspace retained layout tree with header tab and pane children so sub-layer hit-testing uses the same retained layout infrastructure
\item
  settings row children in the settings retained layout so row click resolution uses layout path inspection instead of scroll-offset math
\item
  \texttt{p\allowbreak{}a\allowbreak{}n\allowbreak{}e\allowbreak{}H\allowbreak{}i\allowbreak{}t\allowbreak{}A\allowbreak{}t\allowbreak{}P\allowbreak{}o\allowbreak{}s\allowbreak{}i\allowbreak{}t\allowbreak{}i\allowbreak{}o\allowbreak{}n}, \texttt{s\allowbreak{}e\allowbreak{}t\allowbreak{}t\allowbreak{}i\allowbreak{}n\allowbreak{}g\allowbreak{}s\allowbreak{}R\allowbreak{}o\allowbreak{}w\allowbreak{}A\allowbreak{}t\allowbreak{}P\allowbreak{}o\allowbreak{}s\allowbreak{}i\allowbreak{}t\allowbreak{}i\allowbreak{}o\allowbreak{}n}, and \texttt{i\allowbreak{}s\allowbreak{}I\allowbreak{}n\allowbreak{}s\allowbreak{}i\allowbreak{}d\allowbreak{}e\allowbreak{}S\allowbreak{}e\allowbreak{}t\allowbreak{}t\allowbreak{}i\allowbreak{}n\allowbreak{}g\allowbreak{}s\allowbreak{}D\allowbreak{}r\allowbreak{}a\allowbreak{}w\allowbreak{}e\allowbreak{}r} removed --- all replaced by retained layout path inspection through \texttt{r\allowbreak{}o\allowbreak{}u\allowbreak{}t\allowbreak{}e\allowbreak{}d\allowbreak{}H\allowbreak{}i\allowbreak{}t}
\item
  \texttt{r\allowbreak{}e\allowbreak{}s\allowbreak{}o\allowbreak{}l\allowbreak{}v\allowbreak{}e\allowbreak{}W\allowbreak{}o\allowbreak{}r\allowbreak{}k\allowbreak{}s\allowbreak{}p\allowbreak{}a\allowbreak{}c\allowbreak{}e\allowbreak{}P\allowbreak{}a\allowbreak{}n\allowbreak{}e\allowbreak{}R\allowbreak{}e\allowbreak{}c\allowbreak{}t\allowbreak{}s} extracted as a shared pane geometry resolver for both the workspace layout tree and the mouse handler
\end{itemize}

What the fourth slice lands:

\begin{itemize}
\tightlist
\item
  \texttt{F\allowbreak{}r\allowbreak{}a\allowbreak{}m\allowbreak{}e\allowbreak{}S\allowbreak{}h\allowbreak{}e\allowbreak{}l\allowbreak{}l\allowbreak{}C\allowbreak{}o\allowbreak{}m\allowbreak{}m\allowbreak{}a\allowbreak{}n\allowbreak{}d\allowbreak{}<\allowbreak{}M\allowbreak{}s\allowbreak{}g\allowbreak{}>} discriminated union in \texttt{a\allowbreak{}p\allowbreak{}p\allowbreak{}-\allowbreak{}f\allowbreak{}r\allowbreak{}a\allowbreak{}m\allowbreak{}e\allowbreak{}-\allowbreak{}t\allowbreak{}y\allowbreak{}p\allowbreak{}e\allowbreak{}s\allowbreak{}.\allowbreak{}t\allowbreak{}s} covering every shell mutation and command emission as plain data facts
\item
  a handler table inside \texttt{c\allowbreak{}r\allowbreak{}e\allowbreak{}a\allowbreak{}t\allowbreak{}e\allowbreak{}F\allowbreak{}r\allowbreak{}a\allowbreak{}m\allowbreak{}e\allowbreak{}d\allowbreak{}A\allowbreak{}p\allowbreak{}p} that interprets each command variant through model mutation and TEA command accumulation
\item
  \texttt{d\allowbreak{}r\allowbreak{}a\allowbreak{}i\allowbreak{}n\allowbreak{}S\allowbreak{}h\allowbreak{}e\allowbreak{}l\allowbreak{}l\allowbreak{}C\allowbreak{}o\allowbreak{}m\allowbreak{}m\allowbreak{}a\allowbreak{}n\allowbreak{}d\allowbreak{}B\allowbreak{}u\allowbreak{}f\allowbreak{}f\allowbreak{}e\allowbreak{}r} using \texttt{b\allowbreak{}u\allowbreak{}f\allowbreak{}f\allowbreak{}e\allowbreak{}r\allowbreak{}R\allowbreak{}u\allowbreak{}n\allowbreak{}t\allowbreak{}i\allowbreak{}m\allowbreak{}e\allowbreak{}R\allowbreak{}o\allowbreak{}u\allowbreak{}t\allowbreak{}e\allowbreak{}R\allowbreak{}e\allowbreak{}s\allowbreak{}u\allowbreak{}l\allowbreak{}t} and \texttt{a\allowbreak{}p\allowbreak{}p\allowbreak{}l\allowbreak{}y\allowbreak{}R\allowbreak{}u\allowbreak{}n\allowbreak{}t\allowbreak{}i\allowbreak{}m\allowbreak{}e\allowbreak{}C\allowbreak{}o\allowbreak{}m\allowbreak{}m\allowbreak{}a\allowbreak{}n\allowbreak{}d\allowbreak{}B\allowbreak{}u\allowbreak{}f\allowbreak{}f\allowbreak{}e\allowbreak{}r} to process route results
\item
  mouse routing handler produces shell commands directly instead of delegating to a separate \texttt{h\allowbreak{}a\allowbreak{}n\allowbreak{}d\allowbreak{}l\allowbreak{}e\allowbreak{}F\allowbreak{}r\allowbreak{}a\allowbreak{}m\allowbreak{}e\allowbreak{}M\allowbreak{}o\allowbreak{}u\allowbreak{}s\allowbreak{}e} function
\item
  key routing handler produces shell commands through per-layer helper functions (\texttt{h\allowbreak{}a\allowbreak{}n\allowbreak{}d\allowbreak{}l\allowbreak{}e\allowbreak{}P\allowbreak{}a\allowbreak{}l\allowbreak{}e\allowbreak{}t\allowbreak{}t\allowbreak{}e\allowbreak{}L\allowbreak{}a\allowbreak{}y\allowbreak{}e\allowbreak{}r\allowbreak{}K\allowbreak{}e\allowbreak{}y\allowbreak{}C\allowbreak{}o\allowbreak{}m\allowbreak{}m\allowbreak{}a\allowbreak{}n\allowbreak{}d\allowbreak{}s}, \texttt{h\allowbreak{}a\allowbreak{}n\allowbreak{}d\allowbreak{}l\allowbreak{}e\allowbreak{}S\allowbreak{}e\allowbreak{}t\allowbreak{}t\allowbreak{}i\allowbreak{}n\allowbreak{}g\allowbreak{}s\allowbreak{}L\allowbreak{}a\allowbreak{}y\allowbreak{}e\allowbreak{}r\allowbreak{}K\allowbreak{}e\allowbreak{}y\allowbreak{}C\allowbreak{}o\allowbreak{}m\allowbreak{}m\allowbreak{}a\allowbreak{}n\allowbreak{}d\allowbreak{}s}, \texttt{h\allowbreak{}a\allowbreak{}n\allowbreak{}d\allowbreak{}l\allowbreak{}e\allowbreak{}W\allowbreak{}o\allowbreak{}r\allowbreak{}k\allowbreak{}s\allowbreak{}p\allowbreak{}a\allowbreak{}c\allowbreak{}e\allowbreak{}L\allowbreak{}a\allowbreak{}y\allowbreak{}e\allowbreak{}r\allowbreak{}K\allowbreak{}e\allowbreak{}y\allowbreak{}C\allowbreak{}o\allowbreak{}m\allowbreak{}m\allowbreak{}a\allowbreak{}n\allowbreak{}d\allowbreak{}s}, etc.) instead of a 280-line if-chain in \texttt{u\allowbreak{}p\allowbreak{}d\allowbreak{}a\allowbreak{}t\allowbreak{}e\allowbreak{}(\allowbreak{})}
\item
  \texttt{u\allowbreak{}p\allowbreak{}d\allowbreak{}a\allowbreak{}t\allowbreak{}e\allowbreak{}(\allowbreak{})} key and mouse branches reduced to single-line buffer drains
\item
  removed: \texttt{h\allowbreak{}a\allowbreak{}n\allowbreak{}d\allowbreak{}l\allowbreak{}e\allowbreak{}F\allowbreak{}r\allowbreak{}a\allowbreak{}m\allowbreak{}e\allowbreak{}M\allowbreak{}o\allowbreak{}u\allowbreak{}s\allowbreak{}e}, \texttt{w\allowbreak{}i\allowbreak{}t\allowbreak{}h\allowbreak{}O\allowbreak{}b\allowbreak{}s\allowbreak{}e\allowbreak{}r\allowbreak{}v\allowbreak{}e\allowbreak{}d\allowbreak{}K\allowbreak{}e\allowbreak{}y}, \texttt{a\allowbreak{}p\allowbreak{}p\allowbreak{}l\allowbreak{}y\allowbreak{}Q\allowbreak{}u\allowbreak{}i\allowbreak{}t\allowbreak{}R\allowbreak{}e\allowbreak{}q\allowbreak{}u\allowbreak{}e\allowbreak{}s\allowbreak{}t}, \texttt{o\allowbreak{}b\allowbreak{}s\allowbreak{}e\allowbreak{}r\allowbreak{}v\allowbreak{}e\allowbreak{}d\allowbreak{}R\allowbreak{}o\allowbreak{}u\allowbreak{}t\allowbreak{}e\allowbreak{}F\allowbreak{}o\allowbreak{}r\allowbreak{}L\allowbreak{}a\allowbreak{}y\allowbreak{}e\allowbreak{}r}, \texttt{a\allowbreak{}p\allowbreak{}p\allowbreak{}l\allowbreak{}y\allowbreak{}H\allowbreak{}e\allowbreak{}l\allowbreak{}p\allowbreak{}S\allowbreak{}c\allowbreak{}r\allowbreak{}o\allowbreak{}l\allowbreak{}l\allowbreak{}A\allowbreak{}c\allowbreak{}t\allowbreak{}i\allowbreak{}o\allowbreak{}n}
\end{itemize}

What this cycle does not address:

\begin{itemize}
\tightlist
\item
  notification toast hit-testing stays outside the retained layout system (viewport-positioned overlays managed by notification.ts)
\item
  FrameScopedMsg branches (timers, transitions, notification ticks) stay as special cases in update() --- they are frame-owned state mutations, not routed input events
\end{itemize}

\section{Runtime-buffer-backed command/effect dispatch --- design intent}\label{record-155-runtime-buffer-backed-commandeffect-dispatch-design-intent}

This section captures the architectural intent for the final slice so it is not lost across sessions.

\subsection{Current pattern}\label{record-155-current-pattern}

The shell's \texttt{u\allowbreak{}p\allowbreak{}d\allowbreak{}a\allowbreak{}t\allowbreak{}e\allowbreak{}(\allowbreak{})} function has two separate steps:

\begin{enumerate}
\def\labelenumi{\arabic{enumi}.}
\item
  \textbf{Routing}: \texttt{r\allowbreak{}e\allowbreak{}s\allowbreak{}o\allowbreak{}l\allowbreak{}v\allowbreak{}e\allowbreak{}R\allowbreak{}o\allowbreak{}u\allowbreak{}t\allowbreak{}e\allowbreak{}d\allowbreak{}K\allowbreak{}e\allowbreak{}y\allowbreak{}L\allowbreak{}a\allowbreak{}y\allowbreak{}e\allowbreak{}r\allowbreak{}(\allowbreak{})} / \texttt{r\allowbreak{}e\allowbreak{}s\allowbreak{}o\allowbreak{}l\allowbreak{}v\allowbreak{}e\allowbreak{}R\allowbreak{}o\allowbreak{}u\allowbreak{}t\allowbreak{}e\allowbreak{}d\allowbreak{}M\allowbreak{}o\allowbreak{}u\allowbreak{}s\allowbreak{}e\allowbreak{}L\allowbreak{}a\allowbreak{}y\allowbreak{}e\allowbreak{}r\allowbreak{}(\allowbreak{})} calls \texttt{r\allowbreak{}o\allowbreak{}u\allowbreak{}t\allowbreak{}e\allowbreak{}R\allowbreak{}u\allowbreak{}n\allowbreak{}t\allowbreak{}i\allowbreak{}m\allowbreak{}e\allowbreak{}I\allowbreak{}n\allowbreak{}p\allowbreak{}u\allowbreak{}t\allowbreak{}(\allowbreak{})} to determine which layer owns the input. The routing handler callback returns \texttt{\{\allowbreak{} \allowbreak{}h\allowbreak{}a\allowbreak{}n\allowbreak{}d\allowbreak{}l\allowbreak{}e\allowbreak{}d\allowbreak{}:\allowbreak{} \allowbreak{}t\allowbreak{}r\allowbreak{}u\allowbreak{}e\allowbreak{} \allowbreak{}\}} or \texttt{\{\allowbreak{} \allowbreak{}s\allowbreak{}t\allowbreak{}o\allowbreak{}p\allowbreak{}:\allowbreak{} \allowbreak{}t\allowbreak{}r\allowbreak{}u\allowbreak{}e\allowbreak{} \allowbreak{}\}} --- it never produces commands or effects.
\item
  \textbf{Command production}: After routing, the \texttt{u\allowbreak{}p\allowbreak{}d\allowbreak{}a\allowbreak{}t\allowbreak{}e\allowbreak{}(\allowbreak{})} function matches on \texttt{r\allowbreak{}o\allowbreak{}u\allowbreak{}t\allowbreak{}e\allowbreak{}d\allowbreak{}L\allowbreak{}a\allowbreak{}y\allowbreak{}e\allowbreak{}r\allowbreak{}.\allowbreak{}k\allowbreak{}i\allowbreak{}n\allowbreak{}d} through a long chain of \texttt{i\allowbreak{}f} branches, each producing \texttt{[\allowbreak{}m\allowbreak{}o\allowbreak{}d\allowbreak{}e\allowbreak{}l\allowbreak{},\allowbreak{} \allowbreak{}C\allowbreak{}m\allowbreak{}d\allowbreak{}<\allowbreak{}F\allowbreak{}r\allowbreak{}a\allowbreak{}m\allowbreak{}e\allowbreak{}d\allowbreak{}A\allowbreak{}p\allowbreak{}p\allowbreak{}M\allowbreak{}s\allowbreak{}g\allowbreak{}<\allowbreak{}M\allowbreak{}s\allowbreak{}g\allowbreak{}>\allowbreak{}>\allowbreak{}[\allowbreak{}]\allowbreak{}]} tuples with manually accumulated command arrays.
\end{enumerate}

These are decoupled: routing decides ownership, then entirely separate code produces commands. The runtime buffer API (\texttt{R\allowbreak{}u\allowbreak{}n\allowbreak{}t\allowbreak{}i\allowbreak{}m\allowbreak{}e\allowbreak{}C\allowbreak{}o\allowbreak{}m\allowbreak{}m\allowbreak{}a\allowbreak{}n\allowbreak{}d\allowbreak{}B\allowbreak{}u\allowbreak{}f\allowbreak{}f\allowbreak{}e\allowbreak{}r}, \texttt{R\allowbreak{}u\allowbreak{}n\allowbreak{}t\allowbreak{}i\allowbreak{}m\allowbreak{}e\allowbreak{}E\allowbreak{}f\allowbreak{}f\allowbreak{}e\allowbreak{}c\allowbreak{}t\allowbreak{}B\allowbreak{}u\allowbreak{}f\allowbreak{}f\allowbreak{}e\allowbreak{}r}, \texttt{b\allowbreak{}u\allowbreak{}f\allowbreak{}f\allowbreak{}e\allowbreak{}r\allowbreak{}R\allowbreak{}u\allowbreak{}n\allowbreak{}t\allowbreak{}i\allowbreak{}m\allowbreak{}e\allowbreak{}R\allowbreak{}o\allowbreak{}u\allowbreak{}t\allowbreak{}e\allowbreak{}R\allowbreak{}e\allowbreak{}s\allowbreak{}u\allowbreak{}l\allowbreak{}t}, \texttt{a\allowbreak{}p\allowbreak{}p\allowbreak{}l\allowbreak{}y\allowbreak{}R\allowbreak{}u\allowbreak{}n\allowbreak{}t\allowbreak{}i\allowbreak{}m\allowbreak{}e\allowbreak{}C\allowbreak{}o\allowbreak{}m\allowbreak{}m\allowbreak{}a\allowbreak{}n\allowbreak{}d\allowbreak{}B\allowbreak{}u\allowbreak{}f\allowbreak{}f\allowbreak{}e\allowbreak{}r}) exists but the shell does not use it.

\subsection{Target pattern}\label{record-155-target-pattern}

Input is routed through the runtime view stack. Instead of routing handlers returning bare \texttt{\{\allowbreak{} \allowbreak{}h\allowbreak{}a\allowbreak{}n\allowbreak{}d\allowbreak{}l\allowbreak{}e\allowbreak{}d\allowbreak{}:\allowbreak{} \allowbreak{}t\allowbreak{}r\allowbreak{}u\allowbreak{}e\allowbreak{} \allowbreak{}\}}, they return commands and effects as part of the route outcome:

\begin{CodeBlock}
routeRuntimeInput(stack, layouts, event, ({ layer, hit }) => {
  // Layer-specific command production happens HERE,
  // inside the routing callback
  return {
    handled: true,
    commands: [emitMsgForPage(pageId, action)],
    effects: [],
  };
});
\end{CodeBlock}

The route result carries accumulated commands and effects from all visited layers. The shell then buffers them:

\begin{CodeBlock}
const buffers = bufferRuntimeRouteResult(
  createRuntimeBuffers(),
  routeResult,
);
\end{CodeBlock}

And applies them after routing completes:

\begin{CodeBlock}
const { state, applied } = applyRuntimeCommandBuffer(
  model,
  buffers.commands,
  applyCommand,
);
\end{CodeBlock}

\subsection{Why this matters}\label{record-155-why-this-matters}

\begin{itemize}
\tightlist
\item
  Routing and command production become one step instead of two. The layer that handles the input also says what should happen, inside the same callback.
\item
  The long \texttt{i\allowbreak{}f\allowbreak{} \allowbreak{}(\allowbreak{}r\allowbreak{}o\allowbreak{}u\allowbreak{}t\allowbreak{}e\allowbreak{}d\allowbreak{}L\allowbreak{}a\allowbreak{}y\allowbreak{}e\allowbreak{}r\allowbreak{}.\allowbreak{}k\allowbreak{}i\allowbreak{}n\allowbreak{}d\allowbreak{} \allowbreak{}=\allowbreak{}=\allowbreak{}=\allowbreak{} \allowbreak{}.\allowbreak{}.\allowbreak{}.\allowbreak{})} chain in \texttt{u\allowbreak{}p\allowbreak{}d\allowbreak{}a\allowbreak{}t\allowbreak{}e\allowbreak{}(\allowbreak{})} moves into the routing handler, which already knows the layer context.
\item
  Commands are collected through the same buffer mechanism the runtime engine uses for component-level input handling, instead of through ad-hoc array accumulation.
\item
  Effects become a first-class channel. The shell currently has no explicit effect separation --- everything is a command. The buffer API gives effects their own lane without changing the TEA return type.
\end{itemize}

\subsection{Scope}\label{record-155-scope}

This is a structural refactoring of the shell's update function. It touches \textasciitilde20 functions across \texttt{a\allowbreak{}p\allowbreak{}p\allowbreak{}-\allowbreak{}f\allowbreak{}r\allowbreak{}a\allowbreak{}m\allowbreak{}e\allowbreak{}.\allowbreak{}t\allowbreak{}s}, \texttt{a\allowbreak{}p\allowbreak{}p\allowbreak{}-\allowbreak{}f\allowbreak{}r\allowbreak{}a\allowbreak{}m\allowbreak{}e\allowbreak{}-\allowbreak{}a\allowbreak{}c\allowbreak{}t\allowbreak{}i\allowbreak{}o\allowbreak{}n\allowbreak{}s\allowbreak{}.\allowbreak{}t\allowbreak{}s}, and \texttt{a\allowbreak{}p\allowbreak{}p\allowbreak{}-\allowbreak{}f\allowbreak{}r\allowbreak{}a\allowbreak{}m\allowbreak{}e\allowbreak{}-\allowbreak{}p\allowbreak{}a\allowbreak{}l\allowbreak{}e\allowbreak{}t\allowbreak{}t\allowbreak{}e\allowbreak{}.\allowbreak{}t\allowbreak{}s}. The change is repetitive but mechanical: every \texttt{[\allowbreak{}m\allowbreak{}o\allowbreak{}d\allowbreak{}e\allowbreak{}l\allowbreak{},\allowbreak{} \allowbreak{}c\allowbreak{}m\allowbreak{}d\allowbreak{}s\allowbreak{}]} return site moves into a routing handler that returns \texttt{\{\allowbreak{} \allowbreak{}h\allowbreak{}a\allowbreak{}n\allowbreak{}d\allowbreak{}l\allowbreak{}e\allowbreak{}d\allowbreak{}:\allowbreak{} \allowbreak{}t\allowbreak{}r\allowbreak{}u\allowbreak{}e\allowbreak{},\allowbreak{} \allowbreak{}c\allowbreak{}o\allowbreak{}m\allowbreak{}m\allowbreak{}a\allowbreak{}n\allowbreak{}d\allowbreak{}s\allowbreak{}:\allowbreak{} \allowbreak{}[\allowbreak{}.\allowbreak{}.\allowbreak{}.\allowbreak{}]\allowbreak{} \allowbreak{}\}}.

Outward shell behavior must not change. The same inputs must produce the same model transitions and the same commands.

\subsection{Risks}\label{record-155-risks}

\begin{itemize}
\tightlist
\item
  The routing handler closure becomes large because it absorbs all layer-specific logic. May need to extract per-layer handler functions to keep readability.
\item
  The TEA \texttt{u\allowbreak{}p\allowbreak{}d\allowbreak{}a\allowbreak{}t\allowbreak{}e\allowbreak{}(\allowbreak{})} return type (\texttt{[\allowbreak{}M\allowbreak{}o\allowbreak{}d\allowbreak{}e\allowbreak{}l\allowbreak{},\allowbreak{} \allowbreak{}C\allowbreak{}m\allowbreak{}d\allowbreak{}[\allowbreak{}]\allowbreak{}]}) is consumed by the event bus. The buffer API produces \texttt{R\allowbreak{}u\allowbreak{}n\allowbreak{}t\allowbreak{}i\allowbreak{}m\allowbreak{}e\allowbreak{}B\allowbreak{}u\allowbreak{}f\allowbreak{}f\allowbreak{}e\allowbreak{}r\allowbreak{}s}, which must be drained back into \texttt{C\allowbreak{}m\allowbreak{}d\allowbreak{}[\allowbreak{}]} at the boundary. This is a seam, not a full TEA replacement.
\item
  Key priority logic (\texttt{f\allowbreak{}r\allowbreak{}a\allowbreak{}m\allowbreak{}e\allowbreak{}-\allowbreak{}f\allowbreak{}i\allowbreak{}r\allowbreak{}s\allowbreak{}t} vs \texttt{p\allowbreak{}a\allowbreak{}g\allowbreak{}e\allowbreak{}-\allowbreak{}f\allowbreak{}i\allowbreak{}r\allowbreak{}s\allowbreak{}t}) currently lives outside routing. It needs to move into the routing handler or be resolved before routing begins.
\end{itemize}

